% chapters/cap7_conclusiones/cap7.tex

\chapter{Conclusiones y Trabajo Futuro}
\label{sec:conclusiones}

\section{Conclusiones}

Este trabajo diseñó, implementó y evaluó un algoritmo de \emph{Quantum Smoothed Particle Hydrodynamics} (QSPH) basado en \emph{Continuous-Time Quantum Walks} (CTQW) para la ecuación de advección unidimensional, dando así respuesta al objetivo general planteado en el capítulo~\ref{ch:objetivos}. En el caso de dos partículas, el circuito abierto ---\emph{Matching Decomposition} (sección~\ref{sec:matching_decomposition}), qubit ancilla y medición intermedia con reinicio como modelo de colisión~\cite{Ciccarello_2022} (sección~\ref{sec:implementacion-ancilla}), y hamiltoniano efectivo $J_{\mathrm{eff}}$ (ecuación~\eqref{eq:J-eff}) que compensa el efecto Zeno (sección~\ref{sec:compensacion-zeno})--- reproduce la transferencia y el equilibrio de la discretización SPH clásica~\eqref{eq:sph_advection_discrete}, tanto para $c=10^{-0{,}5}$ (figura~\ref{fig:sim_definitive}) como para $c=10^{-1}$ (figura~\ref{fig:cuantico-c-menor}).

La primera conclusión de diseño es que el CTQW cerrado no basta. El operador $\hat{U}(t)=e^{-iJt}$ (ecuación~\eqref{eq:operador_U_ctqw}) conserva la energía y genera un intercambio periódico de la cantidad advectada ---las oscilaciones sinusoidales de la figura~\ref{fig:resultado_1}--- incompatible con el atractor de la ecuación~\eqref{eq:sph_advection_discrete} en el dímero (sección~\ref{sec:primera_implementacion}). La referencia clásica, en cambio, relaja hacia $u=1/2$ porque en este caso mínimo las dos partículas permanecen fijas dentro del soporte del kernel (ecuación~\eqref{eq:sph_advection_discrete}). La analogía topológica entre el kernel SPH y la matriz de adyacencia sí es válida para construir $H$ (sección~\ref{subsec: analogía}), pero la relajación de la ecuación~\eqref{eq:sph_advection_discrete} solo aparece al abrir el sistema.

Dentro de ese alcance, el esquema presenta ventajas operativas ya discutidas en los capítulos~\ref{cap:implementacion_circuito} y~\ref{cap:resultados}. \emph{Matching Decomposition} evita la descomposición de Pauli y, en dos partículas, reduce cada paso de Trotter a una única rotación $R_x$ (ecuación~\eqref{eq:puerta_rx}), esa $R_x$ es el \emph{matching} trivial del dímero, no la eficiencia reportada por \cite{matching2026}. La lectura compara las poblaciones $P_j$ con los $u_j$ clásicos, como se definió en la sección~\ref{sec:codificacion_estado}. A diferencia del Euler clásico, que es recursivo y exige todos los estados intermedios, el circuito puede concatenar los $r$ bloques y medir únicamente el estado final, como ilustra la figura~\ref{fig:ventaja}.

La validación, no obstante, se limita a $N=2$. No se han realizado ensayos con más partículas. La construcción actual organiza la evolución en bloques de emparejamiento (sección~\ref{sec:matching_decomposition}), de modo que no cubre un número impar de nodos. Tampoco está demostrado que el protocolo abierto escale sin un ancilla y una medida \emph{mid-circuit} por cada par (sección~\ref{sec:implementacion-ancilla}), esa sobrecarga podría erosionar la eficiencia y aún no se ha cuantificado. Por tanto, el algoritmo se considera válido para analizar un sistema de dos partículas SPH junto con las resstricciones impuestas en este trabajo, y no una formulación general de QSPH.

En conjunto, el trabajo aporta una primera implementación CTQW de QSPH, compacta y físicamente coherente en el caso mínimo, pensada como base para un desarrollo futuro.

\section{Líneas de trabajo futuro}

El paso más lógico e inmediato es generalizar el algoritmo. Hay que aprovechar el \emph{matching} (sección~\ref{sec:matching_decomposition}), la codificación de amplitudes y el mapa de Lindblad~\eqref{eq:lindblad} para $N>2$ y comprobar si cada par requiere su propio ancilla y su propia medida intermedia (sección~\ref{sec:implementacion-ancilla}). Esas pruebas dirán si el canal de colisión sigue reproduciendo la relajación de la ecuación~\eqref{eq:sph_advection_discrete} o si el efecto Zeno (sección~\ref{sec:compensacion-zeno}) y el recuento de qubits vuelven el circuito inviable. Como validación futura conviene, además, introducir una métrica cuantitativa (por ejemplo, el error $L^2$ frente a la solución de la ecuación~\eqref{eq:sph_advection_discrete}) y ensayar un caso $N=4$, donde el \emph{matching} deja de ser trivial.

En paralelo conviene anclar la propuesta al hardware real, en el caso de que los bloques se hagan extensos habra que valorar el reducir profundidad, tolerar ruido, tomar en cuenta las conecciones y restricciones de cada computador y medir fidelidad fuera del simulador ideal. El objetivo a medio plazo no es sustituir ya al SPH clásico, sino dejar un QSPH robusto, accesible y, cuando existan procesadores más capaces, competitivo en tiempo con la integración explícita convencional.
